\begin{lemma}
\label{lemma-blowup-module}
Let $R$ be a ring. Let $M$ be a finite $R$-module. Let $f \in R$
be an element such that $M_f$ is finite locally free of rank $r$.
Then there exists a finitely generated ideal $I \subset R$ with
$V(f) = V(I)$ such that for all $a \in I$ with $R' = R[\frac{I}{a}]$
the strict transform
$$
M' = (M \otimes_R R')/a\text{-power torsion}
$$
is locally free of rank $r$.
\end{lemma}

\begin{proof}
Choose a surjection $R^{\oplus n} \to M$. Choose a finite submodule
$K \subset \Ker(R^{\oplus n} \to M)$ such that $R^{\oplus n}/K \to M$
becomes an isomorphism after inverting $f$. This is possible because
$M_f$ is of finite presentation for example by
Algebra, Lemma \ref{algebra-lemma-finite-projective}.
Set $M_1 = R^{\oplus n}/K$
and suppose we can prove the lemma for $M_1$. Say $I \subset R$ is the
corresponding ideal. Then for $a \in I$ the map
$$
M_1' = (M_1 \otimes_R R')/a\text{-power torsion}
\longrightarrow
M' = (M \otimes_R R')/a\text{-power torsion}
$$
is surjective. It is also an isomorphism after inverting $a$ in $R'$
as $R'_a = R_f$, see Algebra, Lemma \ref{algebra-lemma-blowup-in-principal}.
But $a$ is a nonzerodivisor on $M'_1$, whence the displayed map is an
isomorphism. Thus it suffices to prove the lemma in case $M$ is a finitely
presented $R$-module.

\medskip\noindent
Assume $M$ is a finitely presented $R$-module.
Then $J = \text{Fit}_r(M) \subset S$ is a finitely generated ideal.
We claim that $I = fJ$ works.

\medskip\noindent
We first check that $V(f) = V(I)$. The inclusion $V(f) \subset V(I)$ is
clear. Conversely, if $f \not \in \mathfrak p$, then
$\mathfrak p$ is not an element of $V(J)$ by 
Lemma \ref{lemma-fitting-ideal-generate-locally}.
Thus $\mathfrak p \not \in V(fJ) = V(I)$.

\medskip\noindent
Let $a \in I$ and set $R' = R[\frac{I}{a}]$. We may write $a = fb$
for some $b \in J$. By Algebra, Lemmas \ref{algebra-lemma-affine-blowup} and
\ref{algebra-lemma-blowup-add-principal} we see that $J R' = b R'$
and $b$ is a nonzerodivisor in $R'$. 
Let $\mathfrak p' \subset R' = R[\frac{I}{a}]$ be
a prime ideal. Then $JR'_{\mathfrak p'}$ is generated by $b$.
It follows from
Lemma \ref{lemma-principal-fitting-ideal}
that $M'_{\mathfrak p'}$ can be generated by $r$ elements.
Since $M'$ is finite, there exist $m_1, \ldots, m_r \in M'$ and
$g \in R'$, $g \not \in \mathfrak p'$ such that the corresponding map
$(R')^{\oplus r} \to M'$ becomes surjective after inverting $g$.

\medskip\noindent
Finally, consider the ideal $J' = \text{Fit}_{k - 1}(M')$.
Note that $J' R'_g$ is generated by the coefficients of relations between
$m_1, \ldots, m_r$ (compatibility of Fitting ideal with base change).
Thus it suffices to show that $J' = 0$, see
Lemma \ref{lemma-fitting-ideal-finite-locally-free}.
Since $R'_a = R_f$ (Algebra, Lemma \ref{algebra-lemma-blowup-in-principal})
and $M'_a = M_f$ is free of rank $r$ we see that $J'_a = 0$.
Since $a$ is a nonzerodivisor in $R'$ we
conclude that $J' = 0$ and we win.
\end{proof}